\chapter{Introduction}
\label{ch:intro}

\section{Purpose and Scope}

A golf launch monitor reports selected aspects of club delivery, ball launch and subsequent flight.
What can be recovered depends on the device, observation window, calibration and estimation method; no architecture observes every relevant quantity on every shot.
Radar and optical measurements constrain motion in different ways, while impact and aerodynamic models connect quantities that may not all be observed directly.
These models explain the physical relationships; their usefulness does not establish that every commercial product implements the same collision or flight algorithm.

This review is a technical reference on how launch monitors work. It is
written for engineers building or evaluating such a system, for researchers
using one as an instrument, and for anyone who needs to know what a reported
number actually represents. Its goals are:

\begin{enumerate}[itemsep=2pt]
  \item catalogue \emph{publicly documented measurement principles of selected commercial systems}, at the level
        of radar physics, imaging geometry, and signal processing;
  \item establish, parameter by parameter, \emph{what is measured directly,
        what is inferred through a model, and what is essentially
        estimated} --- especially for club-delivery data such as face angle
        and club path;
  \item survey the \emph{patent landscape} as a source of disclosed sensor architectures, signal-processing methods and claim limitations, while distinguishing those disclosures from verified commercial implementations and design-specific rights determinations;
  \item collect the \emph{physics and algorithms} (D-plane, impact mechanics,
        gear effect, aerodynamic coefficients, spectral spin estimation,
        stereo triangulation, Kalman filtering) relevant to different implementation choices; and
  \item translate all of the above into \emph{concrete design guidance},
        organised by sensing architecture rather than by product.
\end{enumerate}

\begin{keypoint}
A note on neutrality. Products are named throughout, and often criticised.
That is because published vendor definitions, patent claims and measured
tolerances are the primary evidence available in this field --- the
peer-reviewed literature is thin, as \cref{ch:accuracy} documents. Naming a
system is a citation, not a recommendation, and no architecture in this
review is presented as the one to build.
\end{keypoint}

\section{The Two Sensing Families}

\textbf{Doppler radar} measures frequency shifts associated with target radial motion.
Calibrated antenna arrays add angular information, but range observability, ambiguity resolution and tracking depend on the waveform, geometry and estimator (\cref{ch:radar}).
A sufficiently complete outdoor track can constrain downrange motion; a truncated indoor track leaves more of the flight to inference.
Even an outdoor result can contain fitted, extrapolated or normalized quantities rather than a directly observed landing point (\cref{ch:flight,ch:params}).
Club-face orientation requires suitable orientation-sensitive observations or additional assumptions; a radial-velocity observation alone does not determine it.

\textbf{Photometric (camera)} systems estimate ball position and rotation from timed images within a finite capture volume.
Camera count, viewing geometry, visible surface features and illumination vary by product (\cref{ch:camera}).
Calibration, exposure timing, occlusion and rotational ambiguity limit the recovered launch and spin estimates; precision must be established for the reported quantity and operating conditions rather than inferred from the sensor label.
Downrange flight outside that volume is calculated from launch estimates and a flight model.
Club-pose recovery likewise needs sufficient visible geometry, such as identifiable markers or trackable head features; overhead placement alone does not establish observability or accuracy.

Hybrid systems combine complementary observations.
For example, TrackMan describes OERT as synchronized camera and radar tracking, whereas Foresight describes GCQuad as a four-camera system~\cite{trackmanoert,foresightgcquad}.
These examples do not support a universal claim that all high-end vendors require both modalities.
Nor does adding a sensor remove calibration or model uncertainty: its value depends on which ambiguity the additional observation resolves and how the observations are synchronized (\cref{ch:survey,ch:accuracy}).

\section{Reading Guide}

\Cref{ch:params} fixes terminology and coordinate conventions
(using explicit vendor definitions where relevant, without assuming universal agreement).
\Cref{ch:impact} develops the impact physics connecting club delivery to
ball launch --- the D-plane model, spin generation, and gear effect --- which
can inform forward simulation or inverse estimation; the implementation and identifiability of each reported quantity must be established separately. \Cref{ch:radar,ch:camera} treat the two sensing families in
depth. \Cref{ch:survey} surveys the commercial devices.
\Cref{ch:patents} compares disclosed measurement mechanisms and their evidentiary limits.
\Cref{ch:flight} covers ball-flight aerodynamics and trajectory estimation.
\Cref{ch:accuracy} reviews the independent validation literature.
\Cref{ch:implications} distills the review into design guidance.

\begin{keypoint}
Throughout, we use a right-handed coordinate system for a right-handed
golfer: $x$ down the target line, $y$ up, $z$ to the golfer's right.
Positive launch azimuth and elevation point rightward and upward, respectively; this convention is not a universal sign rule for every rotation or spin component.
Each club parameter needs its own reference point and event.
For example, TrackMan defines club speed at the head's geometric center just before first contact, while other delivery parameters use a maximum-compression reference~\cite{trackmanclubspeed,trackmanparameterdefinitions}.
Ball-launch kinematics describe the outgoing ball after separation.
\Cref{ch:params} reconciles these conventions before comparisons or impact calculations.
\end{keypoint}
